\par\bigskip\begingroup\small
\makeatletter\tagsleft@true\let\veqno\@@leqno\makeatother
\begin{center}\textls[120]{Erläuterungen zum Abschnitt} [VII.], S. 420—428.\end{center}
\indent Um das Verständnis der im Abschnitt [VII.] zusammengestellten, recht lückenhaften und nicht einheitlich entworfenen Aufzeichnungen zu erleichtern, geben wir hier zunächst eine zusammenfassende Darstellung ihres Inhalts, bei der auch auf die Erläuterung der Einzelheiten des Textes eingegangen werden soll.

\indent Es bedeute wie in der \textit{Theoria motus}, artt. 5., 6. (Werke, VII, 1906, S. 17, 19) \(v\) die wahre, \(E\) die exzentrische, \(M\) die mittlere Anomalie und wie im art. [2.] unseres Textes \(f=\sin\varphi\) die Exzentrizität\textsuperscript{**}); dann gelten die Gleichungen
\begin{equation}\tag*{(I)}
\begin{aligned}M&=E-f\sin E\qquad (\text{\textit{Theoria motus}, Gl. XII, S. 21})\end{aligned}
\end{equation}
\begin{equation}\tag*{(II)}
\operatorname{tang}\frac E2=\sqrt{\frac{1-f}{1+f}}\operatorname{tang}\frac v2
\qquad (\text{\textit{Theoria motus}, Gl. VII, S. 20}),
\end{equation}
\noindent also, siehe die Gl. [6], oben S. 421,
\begin{equation}\tag*{[6]$'$}
\frac{dv}{dM}=\frac{dv}{dE}\cdot\frac{dE}{dM}
=\frac{\sqrt{1-f^2}}{(1-f\cos E)^2}
=\frac{\cos\varphi}{(1-\sin\varphi\,{.}\,\cos E)^2}\mathpunct{\raisebox{.5ex}{.}}
\end{equation}
\par\medskip\noindent\rule{.24\linewidth}{.4pt}\par\smallskip
\noindent\hspace*{\parindent}\textsuperscript{**}) Im art [1.] wird die Exzentrizität mit \(\varepsilon\), in der \textit{Theoria motus} mit \(e\) bezeichnet.
\clearpage\SourceLeaf{446}\ArithmeticHead{439}{R}{BEMERKUNGEN ZUR REIHENLEHRE. ABSCHNITT VII.}
\indent Es sei nun
\begin{equation}\tag*{(III)}
\frac{dv}{dM}=1+C_1\cos M+C_2\cos2M+C_3\cos3M+\cdotp\cdotp\cdotp,
\end{equation}
\noindent woraus sich durch Integration für die sogenannte \textls[120]{Mittelpunktsgleichung} \(v-M\) die Entwicklung
\begin{equation}\tag*{(IV)}
v-M=C_1\sin M+\frac{C_2}{2}\sin2M+\frac{C_3}{3}\sin3M+\cdotp\cdotp\cdotp
\end{equation}
\noindent ergibt. Die \(C_1,C_2,C_3,{.}\;{.}\;{.}\) sind Funktionen der Exzentrizität \(f\); \textls[120]{es handelt sich um eine asymptotische Darstellung von} \(\frac1n C_n\) \textls[120]{für große Werte von} \(n\).

\indent Die Lösung der Gleichung [7], oben S. 421,
\begin{equation}\tag*{[7]$'$}1-\sin\varphi\cos E=0\end{equation}
\noindent ist (siehe die Gln. [8], [11], wo diese Lösung mit \(E\) bezeichnet wird)
\[\mathfrak E=i\log\operatorname{cotang}\frac\varphi2.\]
\noindent Setzt man
\[\mathfrak M=\mathfrak E-f\sin\mathfrak E,\]
\noindent so ist (siehe die Gln. [9], [10], [12], wo diese Größe mit \(M\) bezeichnet wird)
\[\mathfrak M=i\log\operatorname{cotang}\frac\theta2,\]
\noindent für
\[\operatorname{tang}\frac\theta2=\operatorname{tang}\frac\varphi2 e^{\cos\varphi},\]
\noindent und also \(\mathfrak M\) die Lösung der Gleichung
\[1-g\cos M=0,\]
\noindent wo
\[g=\sin\theta\]
\noindent ist. Setzt man nun\textsuperscript{*})
\begin{equation}\tag*{($\ast$)}E=\mathfrak E+\varepsilon,\quad M=\mathfrak M+\mu,\end{equation}
\noindent so ist also
\[\mu=\varepsilon-f\{\sin(\mathfrak E+\varepsilon)-\sin\mathfrak E\}\]
\noindent und, da \(\cos\mathfrak E=\frac1f\) ist (siehe die Gl. [11], oben S. 422),
\begin{equation}\tag*{[13]$'$}
\mu=\frac12 f\sin\mathfrak E\,{.}\,\varepsilon^2+\frac16\varepsilon^3-\frac1{24}f\sin\mathfrak E\,{.}\,\varepsilon^4-\cdotp\;\cdotp,
\end{equation}
\noindent woraus (siehe die Gl. [14], oben S. 422)
\begin{equation}\tag*{[14]$'$}
\frac{d\mu}{d\varepsilon}=1-f\cos E=f\sin\mathfrak E\,{.}\,\varepsilon+\frac12\varepsilon^2-\frac16 f\sin\mathfrak E\,{.}\,\varepsilon^3-\cdotp\cdotp\cdotp
\end{equation}
\noindent folgt. Analog ist dann (siehe die Gl. [15], oben S. 422)
\begin{equation}\tag*{[15]$'$}
1-g\cos M=g\sin\mathfrak M\,{.}\,\mu+\frac12\mu^2-\frac16g\sin\mathfrak M\,{.}\,\mu^3+\cdotp\cdotp\cdotp.
\end{equation}

\indent Wir setzen nun (nach Gl. [30], oben S. 425\textsuperscript{**}))
\par\medskip\noindent\rule{.24\linewidth}{.4pt}\par\smallskip
\noindent\hspace*{\parindent}\textsuperscript{*}) Über die Bedeutung dieser Substitution siehe weiter unten S. 445.
\par\noindent\hspace*{\parindent}\textsuperscript{**}) Die Formeln des art. [2.] von Gl. [16] an sind nichts anderes, als erste Annäherungen der Entwicklungen des art. [4.]; wir verfahren also hier zunächst im Anschluß an diesen art. [4.].
\clearpage\SourceLeaf{447}\ArithmeticHead{440}{L}{ANALYSIS. NACHLASS.}
\[\cos\varphi=p,\quad\cos\theta=q,\quad x=1-g\cos M,\quad y=1-f\cos E,\]
\noindent dann folgt nach [14]$'$
\begin{equation}\tag*{[14]$''$}
\begin{aligned}
\frac1{y^2}={}&\frac{-1}{p^2}\cdot\frac1{\varepsilon^2}\Bigl\{1+\frac1p i\varepsilon+\frac13\varepsilon^2+\frac5{12}\cdot\frac1p i\varepsilon^3+\cdotp\cdotp\cdotp\\
&\hspace{5em}-\frac34\cdot\frac1{p^2}\varepsilon^2-\frac12\cdot\frac1{p^3}i\varepsilon^3-\cdotp\cdotp\cdotp\Bigr\},
\end{aligned}
\end{equation}
\noindent und nach [15]$'$, indem man für \(\mu\) die Reihenentwicklung [13]$'$ einsetzt,
\begin{equation}\tag*{[15]$''$}
\begin{aligned}
x={}&\frac{-pq\varepsilon^2}{2}\Bigl\{1-\frac13\cdot\frac1p i\varepsilon-\frac1{12}\varepsilon^2+\frac1{60}\cdot\frac1p i\varepsilon^3+\cdotp\cdotp\cdotp\\
&\hspace{5em}+\frac14\frac pq\varepsilon^2-\frac16\cdot\frac1q i\varepsilon^3+\cdotp\cdotp\cdotp\Bigr\}.
\end{aligned}
\end{equation}
\noindent Multipliziert man diese beiden Ausdrücke miteinander und mit \(p\), so erhält man zunächst die Gleichung [31], oben S. 425, aus der dann die Gleichung [32] oder
\begin{equation}\tag*{[32]$'$}
\frac{\cos\varphi}{(1-f\cos E)^2}=\frac{\cos\theta}{2(1-g\cos M)}\Bigl\{1-\frac23\operatorname{cotang}\mathfrak E\,{.}\,\varepsilon+\cdotp\cdotp\cdotp\Bigr\}.
\end{equation}
\noindent folgt. Zieht man in [15]$''$ beiderseits die Quadratwurzel aus und berechnet aus der so entstehenden Gleichung \(\varepsilon\) durch Reihenumkehrung, so erhält man die Gleichung [33], oben S. 426, oder
\begin{equation}\tag*{[33]$'$}
\varepsilon=-i\sqrt{\frac{2(1-g\cos M)}{\cos\varphi\,{.}\,\cos\theta}}-\frac13 i\frac{1-g\cos M}{\cos^2\varphi\,{.}\,\cos\theta}-\cdotp\cdotp\cdotp.
\end{equation}
\noindent Setzt man diesen Wert von \(\varepsilon\) in [32] ein, so ergibt sich [34] oder
\begin{equation}\tag*{[34]$'$}
\frac{\cos\varphi}{(1-\sin\varphi\cos E)^2}=\frac{\cos\theta}{2(1-\sin\theta\cos M)}+\frac{\surd2}{3\cos^{\frac32}\varphi}\sqrt{\frac{\cos\theta}{1-\sin\theta\cos M}}+\cdotp\cdotp\cdotp.
\end{equation}

\indent Soweit ist alles ganz streng, und die auftretenden Reihenentwicklungen konvergieren für hinreichend kleine Werte von \(\varepsilon\) und \(\mu\), das heißt also, wenn \(E\) hinreichend nahe an \(\mathfrak E\), und \(M\) hinreichend nahe an \(\mathfrak M\) liegt. Bei der weiteren Rechnung ersetzt nun \textsc{Gauss} die auf den rechten Seiten der Gleichungen [32], [33], [34] auftretenden Reihen durch ihre ersten Glieder; in der Tat sind die Gleichungen [19], [21], [23] des art. [2.], oben S. 422, 423, nichts anderes als die in der angegebenen Weise reduzierten Gleichungen [32]$'$, [33]$'$, [34]$'$. Die Gleichung [23], das heißt [34]$'$ mit Vernachlässigung der nicht hingeschriebenen Glieder, bildet also jetzt den Ausgangspunkt für alles folgende.

\indent Setzt man, wie im art. [3.], oben S. 423, \(\operatorname{tang}\tfrac12\theta=\gamma\), so ergeben die Gleichungen [1]—[5] des art. [1.], oben S. 420\textsuperscript{*}), die Gleichungen [24] und [25], oben S. 423, das heißt die Entwicklungen von
\[\frac{\cos\theta}{1-g\cos M}\quad\text{und}\quad\sqrt{\frac{\cos\theta}{1-g\cos M}}\]
\noindent nach dem Kosinus der Vielfachen von \(M\). In der Entwicklung der Quadratwurzel lautet der Koeffizient von \(\cos nM\)
\[\begin{aligned}
&2(1-\gamma^2)^{\frac12}\frac12\cdot\frac34\cdotp\cdotp\cdotp\frac{2n-1}{2n}\gamma^n\Bigl(1+\frac12\cdot\frac{2n+1}{2n+2}\gamma^2\\
&\hspace{7em}+\frac12\cdot\frac34\cdot\frac{2n+1}{2n+2}\cdot\frac{2n+3}{2n+4}\gamma^4+\cdotp\cdotp\cdotp\Bigr).
\end{aligned}\]
\par\medskip\noindent\rule{.24\linewidth}{.4pt}\par\smallskip
\noindent\hspace*{\parindent}\textsuperscript{*}) Um Übereinstimmung zwischen den Bezeichnungen der artt. [1.] und [3.] zu erzielen, hat man in den Gleichungen [3] und [5] zu nehmen \(\alpha=\gamma=\operatorname{tang}\frac\theta2\), also das dortige \(\theta\) gleich der Hälfte des hier mit \(\theta\) bezeichneten Winkels, und \(u=M\) zu setzen. Man vergl. übrigens den art. [5.] des Abschnitts [III.] über \(F(\alpha,\beta,\gamma,x)\), oben S. 345.
\par\endgroup
\ArithmeticHead{440}{L}{ANALYSIS. NACHLASS.}
